\section{Die Implementierung}
Der Code aus diesem Abschnitt ist in der Datei Basic.lean zu finden.

\subsection{Grundliegende Strukturen}
Meine grundliegenden Strukturen sind Knoten (node), Kanten (edge), die Zeiger der Union-Find-Struktur (unionFindLink) und die Union-Find-Struktur selbst (unionFind).
Knoten habe keine besonderen Eigenschaften, die sie erfüllen müssen, außer, dass sie unterscheidbar sein müssen. Es ist also ausreichend, sie durch eine ID zu repräsentieren.

\begin{lstlisting}
structure (.\node.) : Type where
  mk ::
  id : Nat
\end{lstlisting}

Alternativ hätte ich Knoten auch direkt als natürliche Zahl darstellen können.
Das wäre von der bedeutung her äquivalent gewesen, hätte mir hier und da ein par Zeilen Code gespart und ich hätte keine der Instazen schreiben müssen, weil die für Nat schon gegeben sind.

\begin{lstlisting}
abbrev (.\textcolor{symbolcolor}{node}.) : Type := Nat
\end{lstlisting}

Allerdings wollte ich zum einen meine grundliegenden Strukturen einheitlich halten und zum anderen hatte ich anfänglich noch geplant, den Zeiger Teil des Knoten sein zu lassen.
Davon bin ich aber abgekommen, da ein Graph mit seinen Knoten unabhängig vom Kruskal-Algorithmus existieren,
die Union-Find-Struktur und ihre Zeiger hingegen nur innerhalb des Algorithmus und abhängig von dessen aktuellen Zustand.

Kanten müssen ebenfalls unterscheidbar sein.
Vor allem aber verbinden sie zwei Knoten und haben kosten.
Außerdem sollen Dopplungen von Kanten vermieden werden.
Der Kruskal-Algorithmus ist definiert für ungerichtete Graphen, also sind für Knoten $u$ und $v$, die Kante von $u$ nach $v$ und die Kante von $v$ nach $u$ die slbe Kante.
Um also Dopplungen zu vermeiden, habe ich die Bedingung gestellt, dass die ID des ersten Knoten echt kleiner als die des zweiten sein soll.

\begin{lstlisting}
structure (.\edge.) : Type where
  mk ::
  id : Nat
  node1 : Nat
  node2 : Nat
  cost : Nat
  nodesLt : node1 < node2
\end{lstlisting}

Durch edge.nodesLt habe ich als Nebeneffekt Selbskanten verboten, also Kanten, die einen Knoten mit sich selbst verbinden.
Das habe ich aber bewusst inkauf genommen, da es sich auf das Ergebnis nicht auswirken würde.
Der Algorithmus würde solche Kanten niemals aufnehmen, da jeder Knoten immer in der selben Zusammenhangskomponente ist wie er selbst.
Hätte ich Selbskantennicht von vorneherein ausgeschlossen, hätte ich zur performance steigerung eine Abfrage in den Algorithmus eingebaut,
der diese noch vor dem Prüfen der Zusammenhangskomponente aussortiert.

Dazu sei angemerkt, dass die edge.id eher kosmetischer Natur ist und von mir nur zu debug zwecken benutzt wird.
Es reichen bei korrektem Input edge.node1 und edge.node2 zur eindeutigen identifikation, allerdings finde ich es beim vergleichen übersichtlicher.

Was mir an den Kanten das meiste Kopfzerbrechen bereitet hat, sind die Relationen $=$ und $\le$.
Auf der einen Seite habe ich bestimmt, dass zwei Kanten gleich ($=$) sind, wenn sie in all ihren Atributen überein stimmen.
Das brauchte ich um in Beweisen verschiedene Kanten unterscheiden zu können.
Auf der anderen Seite habe ich festgelegt, dass eine Kante kleinergleich ($\le$) einer andere ist, wenn ihre Kosten kleinergleich sind.
Das war notwendig, um die Kanten-Liste nach Kosten sortieren zu können.
Das Problem, das sich daraus ergibt, ist, dass $\le$ jetzt nicht antisemetrisch ist.
Für zwei Kanten $e$ und $f$ gilt also nicht, dass aus $e \le f \land f \le e$ auch automatisch $e = f$ folgt.
Dadurch konnte ich manche Lemma aus Mathlib über sortierte Listen nicht benutzen, da sie Lineare Ordnung voraus setzen, was die Antisemetrie von $\le$ beinhaltet.

Die Zeiger sind müssen auf den ersten Blick keine Bedingungen erfüllen.
Sie beinhalten nur die Information, zu welchem Knoten sie gehören, auf welchen sie zeigen und ihren Rang.
Allerdings habe ich doch eine Bedingung eingebaut, indem ich den Rang nach oben auf die Länge der Knoten-Liste beschränke.
Damit zeige ich, dass das suchen des Repräsentanten der Zusammenhangskomponenten auch terminiert.

\begin{lstlisting}
structure unionFindLink (nodeList : List node) : Type where
  mk ::
  nodeId : Nat
  ccId : Nat
  rank : Fin nodeList.length
\end{lstlisting}

Die Union-Find-Struktur besteht eigentlich nur aus der Liste an Zeigern, aber für die müssen eine menge Bedingungen gelten.

\begin{itemize}
    \item Für jeden Knoten muss es genau einen Zeiger geben.
    \item Jeder Zeiger zeigt auf einen eindeutigen Knoten.
    \item Es gibt genau so viele Zeiger wie Knoten.
    \item Wenn ein Zeiger nicht auf seinen eigenen Knoten zeigt, dann hat der Zeiger des Knoten auf den er zeigt einen höheren Rang.
    \item Es gibt keine Dopplungen.
    \item Und die Rang-Invariante
\end{itemize}

Die Rang-Invariante besagt eigentlich, dass für einen Zeiger mit Rang $r$ es mindestens $2^r-1$ Zeiger gibt, die direkt oder indirekt auf ihn zeigen.
Allerdings habe ich mich für den Anfang für eine schwächer Version entschieden, die sich einfacher beweisen lässt.
Nämlich, dass es für einen Zeiger mit Rang $r$ mindestens $r$ Zeiger gibt, die nicht auf ihren eigenen Knoten zeigen.

\begin{lstlisting}
structure unionFind (nodeList : List node) : Type where
  mk ::
  linkList : List (unionFindLink nodeList)
  matching_nodeId : ∀ x ∈ nodeList, ∃! y ∈ linkList, x.id = y.nodeId
  matching_ccId : ∀ y ∈ linkList, ∃! x ∈ nodeList, x.id = y.ccId
  matching_length : nodeList.length = linkList.length
  matching_rank (y : unionFindLink nodeList) (h : y ∈ linkList) : y.nodeId = y.ccId ∨ y.rank < (List.choose (fun x => x.nodeId = y.ccId) linkList (exists_parent_link linkList matching_nodeId matching_ccId y h)).rank
  nodup : linkList.Nodup
  rankInvariant (x : unionFindLink nodeList) : x ∈ linkList → (linkList.filter (fun y => y.rank < x.rank ∧ ¬y.nodeId = y.ccId)).length ≥ x.rank
\end{lstlisting}

Um in matching\_rank über den Zeiger des Kontens, auf den der ursprüngliche Zeiger zeigt, aussagen treffen zu können, braucht es den Beweis, das ein solcher auch existiert.
Das folgt direkt aus matching\_nodeId und matching\_ccId, aber zur besseren Übersicht habe ich es in ein eigenes Lemma ausgelagert.

\begin{lstlisting}
theorem exists_parent_link {nodeList : List node} (linkList : List (unionFindLink nodeList)) (matching_nodeId : ∀ x ∈ nodeList, ∃! y ∈ linkList, x.id = y.nodeId) (matching_ccId : ∀ y ∈ linkList, ∃! x ∈ nodeList, x.id = y.ccId) : ∀ x ∈ linkList, ∃ y ∈ linkList, y.nodeId = x.ccId
\end{lstlisting}

\subsection{Initialisierung der Union-Find-Struktur}
Die Funktion, die für eine gültige Knoten-Liste, also eine ohne Duplikate, die Union-Find-Struktur initialisiert, muss zum einen die Zeiger-Liste erstellen
und zum anderen die Nachweise erbringen, dass alle an sie gestellten Bedingungen auch erfüllt werden.

Zuerst habe ich den trivialen Fall einer leeren Knoten-Liste abgefragt.
In diesem Fall ist auch die Zeiger-Liste leer und keiner der Beweise ist komplexer als zwei simp.
Die erstellung der Zeiger-Liste habe ich in eine Verschachtelte Hilfsfunktion ausgelagert, die für jeden Knoten einen Zeiger mit Rang 0 erzeugt, der auf ihn selbst zeigt.
Hierfür war es auch wichtig, erst den trivialen Fall abzufragen, damit jetzt die 0 kleiner als die Länge der Knoten-Liste und damit als Rang zulässig ist.

\begin{lstlisting}
def init_unionFind.linkList (nodeList : List node) (h_nonempty : nodeList.length ≠ 0) : List (unionFindLink (nodeList : List node)) :=
  helper nodeList where
    helper : List node → List (unionFindLink (nodeList : List node))
      | [] => []
      | x::xs => { nodeId := x.id, ccId := x.id, rank := ⟨0, zero_lt_iff.mpr h_nonempty⟩ } :: (helper xs)
\end{lstlisting}

Die Beweise habe ich ebenfalls ausgelagert.
Dadurch wurde die eigentliche Funktion deutlich kleiner und übersichtlicher.

\begin{lstlisting}
def init_unionFind (nodeList : List node) (h_nodup : nodeList.Nodup) : unionFind (nodeList : List node) :=
  if h : nodeList.length ≠ 0
    then
      let linkList : List (unionFindLink nodeList) := (init_unionFind.linkList nodeList h)
      have h_matching_nodeId : ∀ x ∈ nodeList, ∃! y, y ∈ linkList ∧ x.id = y.nodeId := …

      have h_matching_ccId : ∀ y ∈ linkList, ∃! x, x ∈ nodeList ∧ x.id = y.ccId := …

      have h_matching_length : nodeList.length = linkList.length := …

      have h_matching_rank : ∀ (y : unionFindLink nodeList) (h : y ∈ linkList), y.nodeId = y.ccId ∨ y.rank < (List.choose (fun x => x.nodeId = y.ccId) linkList (exists_parent_link linkList h_matching_nodeId h_matching_ccId y h)).rank := …

      have h_nodup : linkList.Nodup := …

      have h_rankInvariant : (z : unionFindLink nodeList) → z ∈ linkList → (linkList.filter (fun y => y.rank < z.rank ∧ ¬y.nodeId = y.ccId)).length ≥ z.rank := …

      ⟨linkList, h_matching_nodeId, h_matching_ccId, h_matching_length, h_matching_rank, h_nodup, h_rankInvariant⟩
    else
      let linkList : List (unionFindLink nodeList) := []
      ⟨linkList, by simp at h; simp [h], by simp [linkList], by simp at h; simp [h, linkList], by intro y h_y_in; simp [linkList] at h_y_in, by simp [linkList], by simp [linkList]⟩
\end{lstlisting}

Um matching\_nodeId zu zeigen, habe ich für einen generischen Knoten den entsprechenden Zeiger gebaut und gezeigt, dass sich dieser auch in der Zeiger-Liste befindet.
Anschließend musste ich zeigen, dass dieser Zeiger auch die Anforderungen erfüllt (was er per definition tut)
und, da ich $\exists!$ gefordert habe, dass jeder Zeiger aus der Liste, der die Anforderungen erfüllt, in Wahrheit der selbe ist.

\begin{lstlisting}
theorem init_unionFind.linkList.matching_nodeId {nodeList : List node} (h_nonempty : nodeList.length ≠ 0) : ∀ x ∈ nodeList, ∃! y, y ∈ (init_unionFind.linkList nodeList h_nonempty) ∧ x.id = y.nodeId := by
  intro x h_in
  set y : unionFindLink nodeList := ⟨x.id, x.id, ⟨0, zero_lt_iff.mpr h_nonempty⟩⟩
  rcases (mem_split h_in) with ⟨l₁, l₂, h_split, h_nin⟩
  have h_y_in : y ∈ (init_unionFind.linkList nodeList h_nonempty) := by
    simp [init_unionFind.linkList, y, h_split, init_unionFind.linkList.helper_append]
    right
    simp [init_unionFind.linkList.helper]
  refine ⟨y, ?_⟩
  constructor
  · constructor
    · exact h_y_in
    · simp [y]
  · intro z h_z
    rcases h_z with ⟨h_z_in, h_z_eq_id⟩
    simp [init_unionFind.linkList] at h_z_in
    have h_z_eq_self := init_unionFind.linkList.helper_all_self_connected nodeList nodeList h_nonempty z h_z_in
    have h_z_eq_cc : x.id = z.ccId := by
      rw [← h_z_eq_self]
      exact h_z_eq_id
    have h_z_rank_zero := init_unionFind.linkList.helper_all_rank_zero nodeList nodeList h_nonempty z h_z_in
    simp [y]
    nth_rewrite 1 [h_z_eq_id]
    rw [h_z_eq_cc, ← h_z_rank_zero]
  \end{lstlisting}

Zuerst habe ich mir dafür jedoch drei Hilfslemma geschrieben, die sich alle durch eine einfache Listen-Induktion lösen ließen.
Nämlich, dass eine geteilte Knoten-Liste eine geteilte Zeiger-Liste bedeutet, dass die Zeiger alle auf ihren eigenen Knoten Zeigen und das sie alle Rang 0 haben.

\begin{lstlisting}
theorem init_unionFind.linkList.helper_append (nodeList l₁ l₂ : List node) (h_nonempty : nodeList.length ≠ 0) : init_unionFind.linkList.helper nodeList h_nonempty (l₁ ++ l₂) = (init_unionFind.linkList.helper nodeList h_nonempty l₁) ++ (init_unionFind.linkList.helper nodeList h_nonempty l₂)

theorem init_unionFind.linkList.helper_all_self_connected (nodeList l : List node) (h_nonempty : nodeList.length ≠ 0) : ∀ x ∈ (init_unionFind.linkList.helper nodeList h_nonempty l), x.nodeId = x.ccId

theorem init_unionFind.linkList.helper_all_rank_zero (nodeList l : List node) (h_nonempty : nodeList.length ≠ 0) : ∀ x ∈ (init_unionFind.linkList.helper nodeList h_nonempty l), x.rank = ⟨0, zero_lt_iff.mpr h_nonempty⟩
\end{lstlisting}

Der Beweis für matching\_ccId ist dem für matching\_nodeId sehr ähnlich nur andersrum, da ich jetzt für einen generischen Zeiger den entsprechenden Knoten baue.
Dem entsprechend brauchte ich auch das Teilungs-Lemma in die andere Richtung.
Das war jedoch komplexer als die Hinrichtung, da ich zuerst definieren musste, wie ich aus einer Zeiger-Liste auf eine Knoten-Liste schließe
und anschließend beweisen, dass es in diesem Fall die selbe ist, aus der die Zeiger-Liste selbst erzeugt wurde.

\begin{lstlisting}
def nodeList_of_unionFindLinkList (nodeList : List node) : List (unionFindLink nodeList) → List node
  | [] => []
  | x::xs => ⟨x.nodeId⟩ :: (nodeList_of_unionFindLinkList nodeList xs)

theorem init_unionFind.linkList.helper_append_reverse (nodeList : List node) (h_nonempty : nodeList.length ≠ 0) (l₁ l₂ : List (unionFindLink nodeList)) (h_split : init_unionFind.linkList.helper nodeList h_nonempty nodeList = l₁ ++ l₂) : ∃ k₁ k₂, nodeList = k₁ ++ k₂ ∧ init_unionFind.linkList.helper nodeList h_nonempty k₁ = l₁ ∧ init_unionFind.linkList.helper nodeList h_nonempty k₂ = l₂ := by
  set k₁ := nodeList_of_unionFindLinkList nodeList l₁
  set k₂ := nodeList_of_unionFindLinkList nodeList l₂
  refine ⟨k₁, k₂, ?_⟩
  have h_eq₂ : init_unionFind.linkList.helper nodeList h_nonempty k₁ = l₁ := …

  have h_eq₃ : init_unionFind.linkList.helper nodeList h_nonempty k₂ = l₂ := …

  have h_eq₁ : nodeList = k₁ ++ k₂ := …

  refine ⟨h_eq₁, h_eq₂, h_eq₃⟩

theorem init_unionFind.linkList.matching_ccId {nodeList : List node} (h_nonempty : nodeList.length ≠ 0) : ∀ y ∈ (init_unionFind.linkList nodeList h_nonempty), ∃! x, x ∈ nodeList ∧ x.id = y.ccId
\end{lstlisting}

Der Beweis für matching\_length lässt sich durch eine einfache Listen-Induktion zeigen, diese hat aber einen Harken.
Da der Typ der Ränge der Zeiger von der Länge der Knoten-Liste abhängt, hätten bei einer Induktion über die Knoten-Liste
die Zeiger-Listen in der Induktionshypothese und im Ziel einen verschiedenen Typ und die Induktionshypothese ließe sich nicht anwenden.
Deshalb führe ich den Beweis nicht über init\_unionFind.linkList und die Knoten-Liste, sonder über die innere Hilfsfunktion und eine beliebige Knoten-Liste.
Gleiches gilt für den Beweis von nodup, den ich durch eine doppelte Induktion gelöst habe.

\begin{lstlisting}
theorem init_unionFind.linkList.helper_matching_length (nodeList : List node) (h_nonempty : nodeList.length ≠ 0) (l : List node) : (init_unionFind.linkList.helper nodeList h_nonempty l).length = l.length

theorem init_unionFind.linkList.helper_nodup (nodeList : List node) (h_nonempty : nodeList.length ≠ 0) (l : List node) (h_nodup : l.Nodup) : (init_unionFind.linkList.helper nodeList h_nonempty l).Nodup
\end{lstlisting}

matching\_rank hat kein eigenes Lemma bekommen, da es quasi durch das Lemma, dass alle Zeiger Rang 0 haben geöst wurde.

\subsection{Find}
Um für einen Zeiger die Zusammenhangskomponente zu finden, geht man die Kette der Zeiger bis zum Ende.
Zeigt der aktuelle Zeiger auf seinen eigenen Knoten ist man am Ende angekommen und der Knoten repräsentiert die Zusammenhangskomponente,
tut er das nicht, geht man weiter zum Zeiger des Knoten, auf den er Zeigt.

Da das eine rekursive Funktion ist, und es kein Element gibt, das in jedem Schritt kleiner wird, braucht Lean einen nachweis, dass diese Funktion auch terminiert.
Über matching\_rank zeige ich, dass der Rang des neuen Zeigers höher ist als der des aktuellen.
Dadurch habe ich doch eine natürlische Zahl, die in jedem Schritt kleiner wird, nämlich die Rang Oberschranke - der Rang des aktuellen Zeigers.
Das ist auch der grund, warum ich den Rang überhaupt beschränke.

\begin{lstlisting}
def connected_component_of_unionFind_of_unionFindLink {nodeList : List node} (uF : unionFind nodeList) (uFL : unionFindLink nodeList) (h_uFL_in : uFL ∈ uF.linkList): Nat :=
  if h_self_connected : uFL.ccId = uFL.nodeId
    then
      uFL.ccId
    else
      let parent : unionFindLink nodeList := List.choose (fun y => y.nodeId = uFL.ccId) uF.linkList (exists_parent_link uF.linkList uF.matching_nodeId uF.matching_ccId uFL h_uFL_in)
      have h_parent_in : parent ∈ uF.linkList := List.choose_mem (fun y => y.nodeId = uFL.ccId) uF.linkList (exists_parent_link uF.linkList uF.matching_nodeId uF.matching_ccId uFL h_uFL_in)
      have h_r : parent.rank > uFL.rank := by
        rcases uF.matching_rank uFL h_uFL_in with h | h
        · simp [h] at h_self_connected
        · have h_parent_eq : parent = List.choose (fun y => y.nodeId = uFL.ccId) uF.linkList (exists_parent_link uF.linkList uF.matching_nodeId uF.matching_ccId uFL h_uFL_in) := by rfl
          rw [← h_parent_eq] at h
          exact h
      connected_component_of_unionFind_of_unionFindLink uF parent h_parent_in
termination_by nodeList.length - uFL.rank
  decreasing_by
  have h : ↑(List.choose (fun y => y.nodeId = uFL.ccId) uF.linkList (exists_parent_link uF.linkList uF.matching_nodeId uF.matching_ccId uFL h_uFL_in)).rank = parent.rank := by
    simp [parent]
  simp [h]
  apply Nat.sub_lt_sub_left
  · exact uFL.rank.isLt
  · exact h_r
\end{lstlisting}

Um die Zusammenhangskomponente für eine gegebene Knoten-ID zu finden, mache ich nun diese Funktion, ich muss nur einmal zeigen, dass es den dazugehörigen Zeiger gibt.

\begin{lstlisting}
def connected_component_of_unionFind_of_id {nodeList : List node} (uF : unionFind nodeList) (id : Nat) (h : ∃ x ∈ nodeList, x.id = id) : Nat :=
  have h_ex_uFL : ∃ uFL ∈ uF.linkList, uFL.nodeId = id := by
    rcases h with ⟨x, h_in, h_eq⟩
    have h_matching_nodeId := uF.matching_nodeId x h_in
    rw [h_eq] at h_matching_nodeId
    rcases h_matching_nodeId with ⟨y, h_in', h_unique⟩
    refine ⟨y, ?_⟩
    simp [h_in']
  let uFL := List.choose (fun x => x.nodeId = id) uF.linkList h_ex_uFL
  have h_uFL_in : uFL ∈ uF.linkList := by
    exact List.choose_mem (fun x => x.nodeId = id) uF.linkList h_ex_uFL
  connected_component_of_unionFind_of_unionFindLink uF uFL h_uFL_in
\end{lstlisting}

\subsection{Union}


\begin{lstlisting}
theorem unionFind.rank_succ_isLt  {nodeList : List node} (uF : unionFind nodeList) (x y : unionFindLink nodeList) : x ≠ y → x ∈ uF.linkList → y ∈ uF.linkList → x.nodeId = x.ccId → y.nodeId = y.ccId → x.rank = y.rank → x.rank.val.succ < nodeList.length
\end{lstlisting}


\begin{lstlisting}

\end{lstlisting}


\begin{lstlisting}

\end{lstlisting}


\begin{lstlisting}

\end{lstlisting}


\begin{lstlisting}

\end{lstlisting}


\begin{lstlisting}

\end{lstlisting}


\begin{lstlisting}

\end{lstlisting}


\begin{lstlisting}

\end{lstlisting}









\subsection{kruskal}
